% s-384-p384-m002; original PDF 384
\PoemStart{s-384-p384-m002}{Künstler und Werk.}
\begin{BookVerse}
\VerseLine{0}{Gelang dir, Mensch, das lang’ erträumte Grosse —}
\VerseLine{1}{Oh das ist Trennung, nimmer zu verein’gen!}
\VerseLine{0}{Die Andern all, Zuschauer sind es, blosse:}
\VerseLine{1}{Du — bist sein Feind, geschieden von dem Dein’gen.}
\Stanza
\VerseLine{0}{Es ist dein Kind: wie solltest du’s nicht pein’gen}
\VerseLine{1}{Mit harter Hand, mit manchem Hieb und Stosse?}
\VerseLine{1}{Es ist dein Werk: wie wolltest du’s nicht rein’gen?}
\VerseLine{0}{Trägt’s doch der Wehen Spur von deinem Schosse.}
\Stanza
\VerseLine{0}{Du hast es froh, so lang’ es wuchs, getragen:}
\VerseLine{1}{Nun sprang’s zu Tage, ward zu Ton und Wort}
\VerseLine{0}{Und nahm des Schaffens Süssigkeit mit fort.}
\Stanza
\VerseLine{0}{Und glaub’, dir ist es besser, zu entsagen:}
\VerseLine{1}{Verbirg dich, dass die That, des Thäters frei,}
\VerseLine{0}{Nicht mehr beschämt, nicht mehr beschattet sei!}
\end{BookVerse}
\PoemEnd
